\begin{abstract}

Lightpanda is a purpose-built headless browser written in Zig that advertises up to 9$\times$ faster execution and 16$\times$ lower memory consumption compared to Google Chrome, positioning itself as a drop-in replacement for Chrome Headless in end-to-end (E2E) test pipelines. This paper presents an empirical evaluation of this claim across three progressively realistic test stages using Playwright and the Chrome DevTools Protocol (CDP).

Stage 1 confirms performance advantages on controlled micro-benchmarks: Lightpanda is approximately 2$\times$ faster and uses 10$\times$ less memory than Chrome against a local web server. Stage 2 reproduces the speedup in a real GitHub Actions CI pipeline running three Playwright tests against TodoMVC, achieving 1.23$\times$ faster total job duration and 1.50$\times$ faster test execution. However, Stage 3 -- a suite of 40 Playwright tests against a realistic Angular-based e-commerce testing application -- reveals a critical incompatibility: 30 out of 40 tests fail on Lightpanda, while all 40 pass on Chrome without modification. Failures stem from missing browser APIs required by Playwright, inability to create isolated browser contexts, absent CDP domain support, incomplete JavaScript engine coverage, and missing locale support.

Furthermore, Lightpanda's published benchmark figures are based on a purpose-built trivial demo page rather than real web applications, raising serious concerns about the ecological validity of their advertised performance claims. We conclude that Lightpanda, in its current beta state, is not a production-ready replacement for Google Chrome Headless in real-world E2E test pipelines.

\end{abstract}
